\section*{Chapitre \ref{ch:iv2:betting-and-kelly} --- Paris et Kelly}
\markboth{\protect\omTocMark{}{Réponses~: Paris et Kelly}}{\protect\omTocMark{}{Réponses~: Paris et Kelly}}

\begin{solution}{iq:iv2:betting-and-kelly:1}
À 4 contre 1, on perd $1/6$ de dollar par dollar misé, et 5 contre 1 est équitable. La marge du bookmaker sur la pièce est de 5~\% de l'argent qu'il encaisse.

Le profit espéré vaut $\tfrac16 \times 4 - \tfrac56 \times 1 = -\tfrac16$. Des cotes équitables le rendent nul, donc $\tfrac16 b = \tfrac56$ et $b = 5$. En général, les cotes équitables sur un événement de probabilité $p$ sont $(1 - p)/p$ contre 1.

Pour la pièce, chaque cote décimale de 1.90 implique une probabilité de $1/1.90 \approx 52.6~\%$, et les deux s'additionnent à environ 105.3~\%. Si 100 sont misés de chaque côté, le bookmaker encaisse 200 et paie $100 \times 1.90 = 190$ quel que soit le côté gagnant. Il garde 10 sur 200, soit 5~\%, sans aucun risque.

\iqwrong{Lire 5.3~\% comme la marge. L'excédent au-delà de 100~\% se mesure sur le paiement. La part de l'argent encaissé vaut $1 - 1/1.0526 = 5~\%$.}
\end{solution}

\begin{solution}{iq:iv2:betting-and-kelly:2}
Misez environ 418.60~\$ sur $X$ et 581.40~\$ sur $Y$. Dans les deux cas, vous récupérez environ 1\,046.51~\$, soit un profit d'environ 46.51~\$, ou 4.65~\%.

Les probabilités implicites valent $1/2.50 = 0.40$ et $1/1.80 \approx 0.5556$, et leur somme vaut $0.9556$, inférieure à 1. Cet écart est l'arbitrage\index{arbitrage}. On mise proportionnellement aux probabilités implicites, de sorte que les deux paiements soient égaux~: $0.40/0.9556 \approx 41.86\%$ sur $X$ et le reste sur $Y$. Chaque dollar rapporte alors $1/0.9556 \approx 1.0465$.

\iqwrong{Répartir l'argent également. Avec 500~\$ sur chaque issue, on obtient 1\,250~\$ si $X$ gagne et 900~\$ si $Y$ gagne, donc on dispose d'un pari rentable sur $X$ et d'aucun gain verrouillé.}
\end{solution}

\begin{solution}{iq:iv2:betting-and-kelly:3}
On mise 10~\% de sa richesse. Le logarithme de la richesse croît alors d'environ 0.0050 par pari, soit environ 0.5~\%.

La richesse après de nombreux paris dépend de la moyenne de son logarithme, si bien qu'on maximise $g(f) = p\ln(1 + f) + q\ln(1 - f)$, avec $q = 1 - p$. En annulant la dérivée, on obtient $f = p - q = 0.10$. C'est le critère de Kelly\index{critère de Kelly}, et, pour un pari à cote égale, il se réduit à l'avantage. Alors $g = 0.55\ln 1.1 + 0.45\ln 0.9 \approx 0.0050$. Une estimation rapide est $\tfrac12 f^2 = 0.005$, soit la moitié du carré de l'avantage.

\iqwrong{Maximiser la richesse espérée, $1 + (p - q)f$. Cette quantité croît avec $f$, elle conseille donc de tout miser, et une seule perte met alors fin au jeu.}
\end{solution}

\begin{solution}{iq:iv2:betting-and-kelly:4}
À 40~\%, on mise 10~\%. À 30~\%, on ne parie pas~: la formule donne $-5~\%$, et le pari inverse n'est pas proposé.

Pour un pari qui rapporte $b$ contre 1, maximiser $p\ln(1 + bf) + q\ln(1 - f)$ donne $f = p - q/b$, soit le profit espéré par dollar, rapporté à la cote. À 40~\%, le profit espéré vaut $0.4 \times 2 - 0.6 = 0.2$, et le diviser par 2 donne 0.1, soit la même valeur que $0.4 - 0.6/2$. À 30~\%, le profit espéré vaut $-0.1$, donc le pari est mauvais.

\iqwrong{«~Dimensionner selon la probabilité de gagner.~» Un outsider à bon prix peut mériter une mise plus forte qu'un pari probable à mauvais prix. L'avantage et la cote décident.}
\end{solution}

\begin{solution}{iq:iv2:betting-and-kelly:5}
La richesse espérée croît de 5~\% par manche. Un chemin typique perd environ 5.1~\% par manche, car le multiplicateur médian sur deux manches vaut $1.5 \times 0.6 = 0.9$. Donc, si l'on doit engager toute sa richesse à chaque manche, il faut s'attendre à perdre, quoi que dise la moyenne.

La moyenne vaut $\tfrac12 \times 1.5 + \tfrac12 \times 0.6 = 1.05$. La croissance d'un chemin est la moyenne des logarithmes, $\tfrac12 \ln 0.9 \approx -0.0527$, donc le multiplicateur typique par manche vaut $\sqrt{0.9} \approx 0.949$. Les deux diffèrent parce que la moyenne est tirée vers le haut par de rares chemins comportant de très nombreux gains, chemins sur lesquels presque personne n'aboutit. Cet écart, environ la moitié de la variance du rendement, est la traînée de volatilité\index{traînée de volatilité}.

\iqwrong{«~Rendement espéré positif, donc tout miser.~» Une hausse de 50~\% suivie d'une baisse de 40~\% laisse la richesse à 0.9 fois sa valeur initiale, quel que soit l'ordre.}
\iqfollow{Que faites-vous de ce pari~?}{On mise une partie de sa richesse. En maximisant $\tfrac12\ln(1 + 0.5f) + \tfrac12\ln(1 - 0.4f)$, on obtient $f = 25~\%$, et la croissance devient positive, $\tfrac12\ln 1.0125 \approx 0.6~\%$ par manche.}
\end{solution}

\begin{solution}{iq:iv2:betting-and-kelly:6}
La richesse finale espérée vaut $1.2^{10} \approx \$6.19$. La probabilité de terminer avec quoi que ce soit vaut $0.6^{10} \approx 0.6~\%$. Le seul chemin où l'on gagne les dix paris rapporte 1\,024~\$, et il fournit à lui seul toute la moyenne. On ne jouerait donc pas ainsi, à moins qu'une seule perte ne soit sans importance.

Chaque pari multiplie la richesse espérée par $0.6 \times 2 = 1.2$, donc maximiser la moyenne conseille bien de tout miser. Elle la maximise en plaçant tout le résultat sur un seul chemin de probabilité 0.6~\%.

\iqwrong{Prendre la moyenne pour ce que l'on peut espérer avoir. Une probabilité de 99.4~\% de ne rien avoir, voilà ce que l'on peut espérer.}
\iqfollow{Et avec la fraction de Kelly de 20~\%~?}{La richesse espérée vaut $1.04^{10} \approx 1.48$, on n'atteint jamais zéro, et même dix pertes laissent $0.8^{10} \approx 0.107$.}
\end{solution}

\begin{solution}{iq:iv2:betting-and-kelly:7}
L'optimum est de 20~\%, avec une croissance d'environ 0.0201 par pari. Sa moitié, soit 10~\%, croît d'environ 0.0150, soit environ 75~\% aussi vite. L'écart type de la variation du logarithme de la richesse baisse d'environ moitié, d'environ 0.199 à environ 0.098.

Au voisinage de l'optimum, la croissance est proche d'une parabole. Un multiple $c$ de la fraction de Kelly conserve une part $c(2 - c)$ de la croissance maximale~; la demi-Kelly\index{Kelly fractionnaire} conserve donc $\tfrac34$ de la croissance pour moitié moins de risque. C'est le compromis que l'on accepte~: le premier quart de croissance auquel on renonce achète la moitié de la volatilité.

L'écart type vaut $\sqrt{pq}\,\ln\tfrac{1 + f}{1 - f}$. Il donne $0.490 \times \ln(1.2/0.8)$ pour le Kelly plein, et $0.490 \times \ln(1.1/0.9)$ pour la moitié.

\iqwrong{Croire que moitié de mise signifie moitié de croissance. La croissance est concave en la fraction misée, si bien que la première moitié de la mise rapporte les trois quarts de la croissance.}
\end{solution}

\begin{solution}{iq:iv2:betting-and-kelly:8}
On mise 50~\% de sa richesse. Le taux de croissance vaut $\tfrac12\ln 1.125 \approx 0.059$ par pari.

On maximise $\tfrac12\ln(1 + f) + \tfrac12\ln(1 - \tfrac12 f)$. La dérivée s'annule lorsque $1 - \tfrac12 f = \tfrac12(1 + f)$, donc $f = \tfrac12$. En général, lorsqu'on gagne une part $b$ de la mise avec la probabilité $p$ et qu'on perd une part $a$ sinon, $f = p/a - q/b$. Ici, cela vaut $\tfrac12/\tfrac12 - \tfrac12/1 = \tfrac12$.

La règle habituelle $p - q/b$ suppose que l'on peut perdre la mise entière, donc $a = 1$. Ici elle donnerait zéro, comme si le pari était un pile ou face équitable à cote égale. La mise de la formule est en réalité le montant que l'on peut perdre~: miser 50~\% met 25~\% de la richesse en jeu.

\iqwrong{Prendre le rendement espéré de 25~\% pour réponse. Une mise dépend de la forme du pari autant que de sa moyenne.}
\end{solution}

\begin{solution}{iq:iv2:betting-and-kelly:9}
On accepte avec une richesse d'au moins 300~\$. Le profit espéré vaut 25~\$, ce qui est positif, mais ne suffit pas, car un investisseur en logarithme paie aussi pour le risque rapporté à la richesse.

On accepte si $\tfrac12\ln(w + 150) + \tfrac12\ln(w - 100) \geq \ln w$, soit $(w + 150)(w - 100) \geq w^2$. Donc $50w - 15\,000 \geq 0$, ce qui donne $w \geq 300$.

Vérification par la formule de Kelly~: il s'agit d'un pari à $b = 1.5$ avec $p = \tfrac12$, donc $f = \tfrac12 - \tfrac12/1.5 = \tfrac16$. Miser 100~\$ correspond à Kelly pour une richesse de 600~\$. C'est le double de Kelly pour 300~\$, le point où la croissance s'annule.

\iqwrong{Dire que tout pari à moyenne positive vaut la peine. Avec une richesse de 110~\$, ce pile ou face risque presque tout pour un faible gain espéré, et l'utilité logarithmique le refuse.}
\end{solution}

\begin{solution}{iq:iv2:betting-and-kelly:10}
Le levier optimal pour la croissance vaut $\mu/\sigma^2 = 0.08/0.0256 = 3.125$. La croissance au-delà du taux d'emprunt vaut $\mu^2/(2\sigma^2) = 12.5~\%$ par an. Avec le ratio de Sharpe $S = \mu/\sigma$, le levier vaut $S/\sigma$, la croissance vaut $S^2/2$, et la volatilité avec levier vaut $S$, ici 50~\%.

Avec un levier $f$, le logarithme de la richesse croît au rythme $g(f) = f\mu - \tfrac12 f^2\sigma^2$, le second terme étant la traînée de volatilité. C'est une parabole dont le sommet se situe en $f^* = \mu/\sigma^2$. En substituant, on obtient $g^* = \tfrac12\mu^2/\sigma^2$.

\iqwrong{Penser que la réponse dépend de l'échelle de la stratégie. Doubler les positions d'une stratégie double $\mu$ et $\sigma$, divise le levier par deux et laisse inchangés $S$ et donc la croissance.}
\iqwhy{C'est ainsi que Kelly apparaît en pratique~: une volatilité cible égale au ratio de Sharpe. Un ratio de Sharpe de 0.5 indique au plus 50~\% de volatilité, et presque personne ne supporte autant.}
\end{solution}

\begin{solution}{iq:iv2:betting-and-kelly:11}
À 40~\%, la croissance vaut $0.6\ln 1.4 + 0.4\ln 0.6 \approx -0.0024$ par pari, donc un chemin typique se contracte, malgré l'avantage positif. La croissance s'annule vers 38.9~\%.

L'optimum est de 20~\%. Résoudre $0.6\ln(1 + f) + 0.4\ln(1 - f) = 0$ numériquement donne $f \approx 0.389$.

La règle empirique vient de la parabole $g(f) \approx (p - q)f - \tfrac12 f^2$, qui est symétrique par rapport à son sommet et possède donc son second zéro en $2f^*$. La vraie fonction contient un logarithme qui décroît plus brutalement lorsque $f$ approche 1, si bien que le côté droit de la courbe est plus raide que le gauche. Le véritable zéro se situe donc sous le double de l'optimum, et miser trop coûte plus que miser trop peu d'une même quantité. À 10~\%, dix points sous l'optimum, la croissance vaut environ 0.0150. À 30~\%, dix points au-dessus, elle vaut environ 0.0147.

\iqwrong{«~Plus d'avantage, plus de mise, plus de croissance.~» Au-delà de l'optimum, on paie davantage en traînée de volatilité qu'on ne gagne en avantage, et à peu près au double on ne gagne plus rien.}
\iqwhy{Un avantage estimé deux fois supérieur à la réalité transforme une stratégie qui devrait capitaliser en une stratégie qui ne va nulle part, avec tout le risque.}
\end{solution}

\begin{solution}{iq:iv2:betting-and-kelly:12}
On mise 25~\%. Le taux de croissance vaut $0.3\ln 1.125 \approx 0.0353$ par pari. L'issue de remboursement intégral n'ajoute rien au logarithme~: elle ne joue donc aucun rôle dans la taille, compte tenu des deux autres probabilités.

On maximise $0.3\ln(1 + 2f) + 0.4\ln 1 + 0.3\ln(1 - f)$. La dérivée s'annule lorsque $0.6/(1 + 2f) = 0.3/(1 - f)$, donc $0.6 - 0.6f = 0.3 + 0.6f$, et $f = \tfrac14$. Alors $0.3\ln 1.5 + 0.3\ln 0.75 = 0.3\ln 1.125$.

Considérons le même pari sans remboursements intégraux, où un gain ou une perte équiprobable paie 2 contre 1. On obtient là aussi $f = \tfrac12 - \tfrac12/2 = 25~\%$. Le rapport entre les deux probabilités fixe la taille. La croissance est proportionnelle à la part des paris qui se règlent~: elle vaut $0.6$ fois celle de ce pari, soit $0.5\ln 1.125 \approx 0.0589$.

\iqwrong{Abaisser la probabilité de gain à 30~\% dans la formule à deux issues, ce qui donne $0.3 - 0.7/2 < 0$, et donc aucun pari. La formule avec une perte totale ne s'applique que lorsqu'il n'y a que deux issues.}
\end{solution}

\begin{solution}{iq:iv2:betting-and-kelly:13}
Miser environ 19.2~\% sur chacun, soit environ 38.5~\% au total. La croissance par manche vaut environ 0.0395, contre environ 0.0403 pour faire les deux paris l'un après l'autre à 20~\%.

Si les deux paris gagnent, on dispose de $1 + 2f$, si un seul gagne, de 1, et si les deux perdent, de $1 - 2f$. Les probabilités valent 0.36, 0.48 et 0.16. La dérivée de $0.36\ln(1 + 2f) + 0.16\ln(1 - 2f)$ s'annule en $0.72(1 - 2f) = 0.32(1 + 2f)$, donc $f = 0.4/2.08 \approx 0.1923$.

On mise moins de 20~\% chacun parce que le second pari ne peut pas utiliser le résultat du premier. Les mises sont prélevées sur la même richesse et sont exposées d'un seul coup. La perte combinée de 38.5~\% est donc un coup plus dur que deux pertes successives à 20~\%, qui laissent 64~\%.

\iqwrong{Additionner les deux fractions de 20~\%. Pour de petits avantages c'est proche, mais c'est faux en principe, et l'erreur croît avec le nombre de paris simultanés.}
\iqfollow{Et avec beaucoup de paris de ce type~?}{La mise totale augmente, mais reste inférieure à la somme des fractions isolées. Avec 5 paris de ce type, l'optimum est d'environ 16.5~\% sur chacun, soit environ 83~\% au total. Avec 50, il est d'environ 2~\% sur chacun, juste en deçà du seuil où les perdre tous vous ramènerait à zéro.}
\end{solution}

\begin{solution}{iq:iv2:betting-and-kelly:14}
Pour une corrélation de 0.5, chacun reçoit environ 1.67 de levier, soit 3.33 au total, et la croissance vaut environ 16.7~\%. Pour une corrélation nulle, chacun reçoit 2.5, soit 5 au total, et la croissance vaut 25~\%.

Pour plusieurs actifs, $f^* = \Sigma^{-1}\mu$, et $g^* = \tfrac12\mu^\top\Sigma^{-1}\mu$. Par symétrie les deux pondérations sont égales, si bien que chaque pondération est solution de $\sigma^2(1 + \rho)f = \mu$. Cela donne $0.04 \times 1.5 \times f = 0.1$, donc $f = 1.667$. Pour $\rho = 0$, $0.04f = 0.1$ donne $f = 2.5$.

La croissance est égale à la moitié du carré du ratio de Sharpe du meilleur mélange, $\tfrac12 \times 2 \times 0.1 \times f$. Soit $0.1667$ pour $\rho = 0.5$, et $0.25$ pour $\rho = 0$. Le ratio de Sharpe du mélange vaut $0.5\sqrt{2/(1 + \rho)}$, soit environ 0.577 pour 0.5 et environ 0.707 pour 0.

\iqwrong{Additionner les deux leviers de Kelly isolés de 2.5 chacun, ce qui ignore la corrélation. Avec une corrélation de 0.5, un tel levier porte plus de volatilité que ne l'exige le mélange, et la croissance tombe d'environ 16.7~\% à 12.5~\%.}
\end{solution}

\begin{solution}{iq:iv2:betting-and-kelly:15}
Miser 50~\% sur le tuyau à chaque fois. La croissance vaut $\ln 2 - H(0.75) \approx 0.131$ par lancer, où $H$ est l'entropie de la justesse du tuyau, en unités naturelles. Avec un informateur parfait elle vaut $\ln 2 \approx 0.693$, et avec un informateur qui a raison une fois sur deux elle est nulle.

Du côté du parieur, le pari est à la cote 1 contre 1, gagné avec une probabilité de 75~\%, donc $f = 0.75 - 0.25 = 0.5$. La croissance vaut $0.75\ln 1.5 + 0.25\ln 0.5$. On la réécrit sous la forme $\ln 2 + 0.75\ln 0.75 + 0.25\ln 0.25$, qui est la formule. L'article de Kelly de 1956 établit ce lien~: le taux de croissance égale l'information que le tuyau apporte sur le lancer.

Un informateur parfait permet de tout miser et de doubler à chaque fois. Un informateur inutile ne donne aucune raison de parier, la croissance est donc nulle, et à 25~\% de justesse, on parie contre le tuyau et on obtient les mêmes 0.131.

\iqwrong{Dire qu'un informateur à 75~\% donne 75~\% de la croissance parfaite. Il en donne environ 19~\%, car l'information chute vite à mesure que la justesse s'éloigne de 100~\%.}
\iqwhy{Elle donne un prix à un signal. Un signal vaut ce qu'il ajoute à la croissance, et un petit gain de justesse près du pile ou face n'ajoute guère.}
\end{solution}

\begin{solution}{iq:iv2:betting-and-kelly:16}
On répartit la richesse 50-30-20, proportionnellement aux probabilités de gain. La croissance vaut environ 0.069 par course. La répartition ne change pas avec les gains versés.

Avec une répartition $w$, la croissance vaut $\sum_i p_i\ln(w_i o_i) = \sum_i p_i\ln o_i + \sum_i p_i\ln w_i$. Les gains n'apparaissent que dans le premier terme, et les $w_i$ seulement dans le second. Ce second terme est maximal en $w_i = p_i$ lorsque les poids doivent faire 1, fait que l'on appelle inégalité de Gibbs. C'est le pari proportionnel\index{pari proportionnel}. À 3 pour 1, la croissance vaut $\sum_i p_i\ln(3p_i) = \ln 3 - H(p) \approx 1.0986 - 1.0297 = 0.069$.

Avec des gains de 2, 4 et 8, la répartition reste 50-30-20, et la croissance devient $0.5\ln 1 + 0.3\ln 1.2 + 0.2\ln 1.6 \approx 0.149$.

Le résultat repose sur l'obligation de tout miser. Lorsqu'on peut en garder une partie, et que les cotes du bookmaker impliquent des probabilités dont la somme dépasse 1, on peut garder du cash et écarter entièrement un mauvais cheval.

\iqwrong{Tout miser sur le favori. C'est le cheval le plus probable, mais une seule défaite ruine tout, et la croissance vaut $-\infty$.}
\end{solution}

\begin{solution}{iq:iv2:betting-and-kelly:17}
À 100~\$, un pari de 10~\$ représente 10~\% de la richesse, et la croissance vaut environ $-0.0030$ par pari~: il ne faut donc pas le prendre. La croissance s'annule à une richesse d'environ 250~\$. Une mise de 10~\$ est la mise de Kelly pour une richesse de 500~\$.

La mise de Kelly vaut $0.51 - 0.49 = 2~\%$. Ensuite $0.51\ln 1.1 + 0.49\ln 0.9 \approx -0.0030$. La croissance s'annule lorsque $f \approx 0.0400$, soit à peu près le double de Kelly, si bien que la richesse vaut $\$10/0.04 \approx \$250$. Et $\$10/0.02 = \$500$.

Ainsi, avec un petit avantage et une mise minimale, il existe une richesse en dessous de laquelle un bon pari devient mauvais. Entre 250~\$ et 500~\$, un pari de 10~\$ est surdimensionné mais croît encore, et il peut valoir la peine d'être pris. En dessous de 250~\$, le prendre fait perdre, quelle que soit l'allure de son espérance.

\iqwrong{«~L'avantage est positif, donc on joue.~» L'espérance de la richesse augmente de 20 cents par pari, mais la croissance est négative~: c'est le signe que la richesse est trop faible pour prendre ce pari.}
\end{solution}

\begin{solution}{iq:iv2:betting-and-kelly:18}
Le gain espéré est de 10~\$ par pari, et la probabilité d'être un jour ruiné vaut $(2/3)^{20} \approx 0.03~\%$. Une fois la richesse au-delà de 250~\$, Kelly demanderait une mise supérieure à la limite~: la richesse croît donc linéairement, d'environ 10~\$ par pari, tandis que la stratégie de Kelly croîtrait de manière composée.

Avec des mises fixées à 50~\$, il s'agit de la ruine du joueur\index{ruine du joueur} avec 20 pas de capital face à un adversaire infiniment riche. La probabilité de perdre un jour 20 pas nets vaut $(q/p)^{20}$, avec $q/p = \tfrac23$. Le gain espéré vaut $50 \times (0.6 - 0.4)$.

La limite change ce qui est optimal. Avec une richesse bien supérieure à 250~\$, on mise toujours la limite. Or, la croissance en logarithme diminue à mesure que la richesse augmente, et le risque de ruine est minuscule~: on aurait donc intérêt à disposer de plus de marge, ce qui est une raison de chercher davantage de tables, et non de modifier sa mise.

\iqwrong{Miser 20~\% de sa richesse jusqu'à atteindre la limite, en croyant que la formule de Kelly tient encore au-delà. Au-delà de ce point, il ne reste plus de décision, seulement le nombre de paris.}
\iqfollow{Même limite, avec 200~\$~?}{À 200~\$, la mise de 50~\$ représente 25~\%, au-dessus de Kelly mais en deçà du double. La probabilité de ruine vaut $(2/3)^4 \approx 20~\%$ si l'on mise toujours la limite, donc il faut miser moins qu'elle tant que la richesse n'a pas augmenté.}
\end{solution}

\begin{solution}{iq:iv2:betting-and-kelly:19}
À 40~\%, il faut tout miser d'un coup, pour une chance de 40~\%. Des paris de 1~\$ donnent environ $2.5 \times 10^{-18}$. À 60~\%, il faut faire de petits paris, pour une chance à environ $2.5 \times 10^{-18}$ près de 1, alors qu'un unique pari tout-ou-rien donne 60~\%.

Chaque cas découle de la ruine du joueur~: avec $r = q/p$, la probabilité d'atteindre 200 depuis 100 par pas de 1~\$ vaut $(1 - r^{100})/(1 - r^{200}) = 1/(1 + r^{100})$. Pour $p = 0.4$, $r = 1.5$, et $1.5^{-100} \approx 2.5 \times 10^{-18}$. Pour $p = 0.6$, $r = \tfrac23$, et la probabilité vaut $1 - (\tfrac23)^{100}$, à peu de chose près.

Dans un jeu défavorable, chaque pari coûte en moyenne, donc on en veut le moins possible~: c'est le jeu audacieux\index{jeu audacieux}. Dans un jeu favorable, chaque pari rapporte en moyenne, et de petites mises laissent la loi des grands nombres faire le travail.

\iqwrong{«~Kelly est optimal, donc il faut miser Kelly.~» Kelly maximise la croissance à long terme, et non la probabilité d'atteindre un objectif. À 40~\%, elle conseille de ne pas miser, et l'on rate alors à coup sûr l'objectif.}
\iqwhy{Un objectif modifie le pari. Un desk qui doit atteindre un objectif en fin d'année, sans rien à perdre, a la même raison que ce joueur de relever son risque, et c'est pourquoi les limites et les structures de rémunération comptent.}
\end{solution}

\begin{solution}{iq:iv2:betting-and-kelly:20}
La perte espérée est de 50~\$. La prime maximale qu'on paierait pour 20\,000~\$ est d'environ 57.45~\$, donc à 60~\$ on n'achète pas. Pour 10\,000~\$, la même formule donne environ 69.08~\$, donc on achète.

La prime $\pi$ rend indifférent lorsque
\[
\ln(W - \pi) = 0.99\ln W + 0.01\ln(W - 5000),
\]
d'où $\pi = W\left(1 - (1 - 5000/W)^{0.01}\right)$. Pour $W = 20\,000$, cela donne $20\,000 \times (1 - 0.75^{0.01}) \approx 57.45$. Pour $W = 10\,000$, on obtient $10\,000 \times (1 - 0.5^{0.01}) \approx 69.08$.

À 20\,000~\$, on paierait jusqu'à 7.45~\$ au-dessus de la perte espérée, et cette prime demande 10~\$ de plus. Le même risque pèse davantage sur une richesse plus petite~: à 10\,000~\$, on paierait jusqu'à 19.08~\$ au-dessus. Une personne bien plus riche refuserait toute prime nettement supérieure à 50~\$.

\iqwrong{Refuser parce que la perte espérée est de 50~\$, donc toute prime au-dessus est un mauvais pari. C'est vrai en espérance, et pourtant les personnes disposant d'un patrimoine limité doivent parfois payer davantage que la perte espérée.}
\end{solution}

\begin{solution}{iq:iv2:betting-and-kelly:21}
Le gain moyen est de 5\,000~\$, l'écart type de 1\,500~\$, et la probabilité de perte d'environ 0.04~\%. Un investisseur logarithmique accepte un lancer à partir de 200~\$ de richesse. Le collègue n'est pas cohérent si ce refus tient pour chaque richesse que le collègue pourrait atteindre.

Chaque lancer a une moyenne de 50 et un écart type de 150. Sur 100 lancers, la moyenne est multipliée par 100 et l'écart type par 10. On ne perd que si l'on obtient au plus 33 piles, car $200h - 100(100 - h) < 0$ exige $h < 33.3$. Cette probabilité vaut environ 0.00044. Pour un seul lancer, on accepte si $(w + 200)(w - 100) \geq w^2$, soit $w \geq 200$.

C'est le sophisme de Samuelson\index{sophisme de Samuelson}. Si une personne aux préférences fixes refuse un lancer à chaque richesse qui peut survenir dans la séquence, alors elle doit refuser toute la séquence, par récurrence à rebours depuis le dernier lancer. Refuser un lancer et en accepter beaucoup ne peut être cohérent que si le refus vaut pour certaines richesses et non pour d'autres. Un investisseur logarithmique n'en fournit pas d'exemple~: en dessous de 200~\$, il refuse un lancer, et pour toute richesse jusqu'à 10\,000~\$ il refuse aussi l'engagement sur 100 lancers, car 100 faces le ruineraient. Au-delà de 10\,000~\$, il accepte les deux.

\iqwrong{«~La diversification dans le temps rend le pari sûr.~» Les enjeux s'additionnent, donc l'écart type du gain croît lui aussi, d'un facteur 10. La probabilité de perte baisse, mais la taille possible de la perte grandit jusqu'à 10\,000~\$.}
\iqwhy{Cet argument revient dans le débat sur la question de savoir si les longs horizons justifient davantage d'actions. La réponse dépend de l'utilité, et pas seulement de l'arithmétique de la moyenne.}
\end{solution}

\begin{solution}{iq:iv2:betting-and-kelly:22}
Au Kelly plein, le logarithme de la richesse croît de 0.5 par an, une trajectoire typique double en environ 1.39 ans, et la volatilité vaut 100~\%. Au demi-Kelly, il croît de 0.375 par an, double en environ 1.85 ans, et a une volatilité de 50~\%.

Le levier de Kelly rend la volatilité de la stratégie égale à $S$, soit ici 100~\%. Un multiple $c$ donne une volatilité $cS$ et une croissance $(c - \tfrac12 c^2)S^2$. Pour $c = 1$, la croissance vaut 0.5, et pour $c = \tfrac12$ elle vaut 0.375. Le temps de doublement vaut $\ln 2/g$.

Le demi-Kelly met donc un tiers de temps en plus pour doubler, soit environ 0.46 an de plus, pour une volatilité moitié moindre. Un desk jugerait cela une bonne affaire, et c'est pourquoi presque personne n'applique réellement le Kelly plein.

\iqwrong{Utiliser $S$ comme taux de croissance, ce qui donnerait un temps de doublement de 0.69 an. La croissance vaut $S^2/2$, car le frein de volatilité absorbe la moitié du rendement excédentaire au Kelly plein.}
\end{solution}

\begin{solution}{iq:iv2:betting-and-kelly:23}
Au levier 2, la croissance vaut 10.88~\% par an, soit environ 87~\% du maximum de 12.5~\%. La croissance revient à zéro au levier 6.25, soit le double de l'optimum.

La croissance vaut $g(f) = 0.08f - \tfrac12 \times 0.0256f^2 = 0.08f - 0.0128f^2$. Pour $f = 2$~: $0.16 - 0.0512 = 0.1088$. Exprimée en multiple $c = 2/3.125 = 0.64$ du Kelly, la part vaut $c(2 - c) = 0.64 \times 1.36 \approx 0.87$. Et $g = 0$ pour $f = 0.08/0.0128 = 6.25$.

Un cap nettement inférieur à l'optimum ne coûte donc presque rien. La limite prend 64~\% du risque et conserve 87~\% de la croissance.

\iqwrong{Prendre 3.125 comme levier à appliquer parce que c'est le sommet. Dans ce cadre gaussien, le sommet est plat, si bien que s'en écarter vers le bas coûte très peu. En revanche, le risque de surestimer $\mu$ est grand.}
\end{solution}

\begin{solution}{iq:iv2:betting-and-kelly:24}
L'espérance du gain est infinie. Avec un cap de $2^{20}$, elle vaut 21~\$. Un investisseur logarithmique disposant de 1\,000~\$ paierait environ 10.95~\$.

Chaque $k$ contribue $2^{-k} \times 2^k = 1$, donc la somme diverge~: c'est le paradoxe de Saint-Pétersbourg\index{paradoxe de Saint-Pétersbourg}. Avec le plafond, les vingt premiers termes totalisent 20, et toutes les séquences plus longues réunies ont une probabilité $2^{-20}$ et paient $2^{20}$, ce qui ajoute encore 1.

Pour une utilité logarithmique, résoudre $\sum_{k \geq 1}2^{-k}\ln(W - c + 2^k) = \ln W$ en $c$. Pour $W = 1\,000$, cela donne environ 10.95. Ce prix croît lentement avec la richesse, pour atteindre environ 20.87 à 1\,000\,000~\$.

\iqwrong{Répondre «~l'infini, donc n'importe quel prix~». Doubler la fortune de la maison n'ajoute que 1~\$ à la fair value, de sorte que le prix est à peu près le logarithme de ce que la maison peut réellement payer.}
\end{solution}

\begin{solution}{iq:iv2:betting-and-kelly:25}
L'erreur type vaut environ 0.5. La croissance au Kelly plein estimé vaut 0.375 par an si le ratio de Sharpe réel vaut 1, et $-0.375$ s'il vaut 0.5. Au demi-levier de Kelly estimé, elle vaut environ 0.469 et environ 0.094.

Pour des rendements à peu près gaussiens, l'erreur type du ratio de Sharpe, en termes annuels, vaut environ $1/\sqrt{T}$, ici $1/\sqrt 4$. Dimensionné sur l'estimation $\hat S$, le portefeuille a une volatilité égale à $\hat S$. Si le ratio de Sharpe réel vaut $S$, la croissance annuelle vaut $\hat S S - \tfrac12\hat S^2$.

Avec $\hat S = 1.5$, la croissance vaut $1.5S - 1.125$. Ainsi $S = 1$ donne 0.375, et $S = 0.5$ donne $-0.375$. À demi-taille, la volatilité du portefeuille vaut 0.75, et la croissance vaut $0.75S - 0.28125$. Ainsi $S = 1$ donne 0.469, et $S = 0.5$ donne 0.094.

Une estimation de 1.5 avec une erreur type de 0.5 est tout à fait compatible avec une vraie valeur de 0.5. Le Kelly plein sur cette valeur fait perdre de l'argent, alors que le demi-Kelly continue de faire croître la richesse. Le demi-Kelly ne sacrifie pas grand-chose si la vraie valeur est 1.5~: 0.84 contre 1.125 par an. C'est l'argument en faveur des fractions de Kelly sur des avantages estimés.

\iqwrong{Y voir une simple affaire de goût pour le risque. Cela tient même pour une personne qui veut une croissance pure, car un ratio de Sharpe estimé est bruité et souvent trop élevé.}
\iqfollow{Est-ce que 1.5 est le bon centre de l'estimation~?}{Si l'on a choisi le meilleur de nombreuses stratégies, l'estimation est biaisée à la hausse. Une règle de dimensionnement doit donc commencer par la réduire, et seulement ensuite appliquer une fraction.}
\end{solution}

\begin{solution}{iq:iv2:betting-and-kelly:26}
Au Kelly plein, $1/W_t$ est une martingale, ce qui majore par $x$ la probabilité de jamais tomber à $x$. En temps continu, la probabilité au multiple $c$ vaut $x^{2/c - 1}$~: soit $x$ au Kelly plein et $x^3$ au demi-Kelly. Ainsi la probabilité de jamais perdre la moitié vaut 50~\% au Kelly plein et 12.5~\% au demi-Kelly.

Pour la borne discrète, avec une mise $f = 0.2$~:
\[
E\left[\frac{1}{1 + fX}\right] = \frac{0.6}{1.2} + \frac{0.4}{0.8} = 1,
\]
donc $M_t = 1/W_t$ est une martingale positive qui part de 1. Pour une telle martingale, la probabilité d'atteindre un jour le niveau $1/x$ est au plus $x$, par l'inégalité maximale\index{inégalité maximale de Doob} (ou par arrêt optionnel). C'est la borne. Elle n'est pas une égalité, car $W$ saute par-dessus le niveau~: un calcul exact sur le réseau donne environ 45.5~\% pour $x = \tfrac12$, contre 50~\%. Au demi-Kelly discret, elle donne environ 10.5~\%.

En temps continu, $\ln W$ est un mouvement brownien avec dérive $m = (c - \tfrac12 c^2)S^2$ et de volatilité $v = cS$. La probabilité de descendre un jour de $L$ vaut $e^{-2mL/v^2}$, ce qui égale $x^{2m/v^2}$ pour $L = -\ln x$. Ici $2m/v^2 = (2 - c)/c = 2/c - 1$.

\iqwrong{«~La stratégie de Kelly ne tombe jamais à zéro, elle est donc sûre.~» Elle n'atteint jamais zéro, et pourtant il y a une chance de 50~\% de subir un drawdown de moitié et une chance de 10~\% d'un drawdown de 90~\%, quel que soit le point de départ. La plupart des gens ne peuvent pas vivre avec cela.}
\end{solution}

\begin{solution}{iq:iv2:betting-and-kelly:27}
Le multiple le plus élevé est $c \approx 0.376$, et il conserve environ 61~\% de la croissance maximale.

La probabilité de tomber un jour à une fraction $x$ de sa valeur initiale vaut $x^{2/c - 1}$. Fixons-la à 0.05 pour $x = \tfrac12$~: $2/c - 1 = \ln 0.05/\ln 0.5 \approx 4.32$, donc $c = 2/5.32 \approx 0.376$. La part de la croissance maximale conservée vaut $c(2 - c) = 0.376 \times 1.624 \approx 0.61$.

La dérivation découle de la dérive et de la volatilité du logarithme de la richesse. Avec une dérive $m = (c - \tfrac12 c^2)S^2$ et une volatilité $v = cS$, la probabilité que le mouvement brownien avec dérive descende un jour de $L$ sous son point de départ vaut $e^{-2mL/v^2}$. Elle ne dépend pas de $S$ du tout~: seul compte le multiple de Kelly, ce qui mérite d'être souligné.

Une limite sur le drawdown s'exprime donc en unités de Kelly, et les investisseurs qui refusent de perdre la moitié demandent en réalité une fraction d'environ 0.38 de Kelly.

\iqwrong{Essayer de le faire avec une valeur en risque (VaR) quotidienne. Une limite sur un drawdown sur toute la durée exige la probabilité d'atteinte de la trajectoire, et non une distribution à un jour.}
\iqfollow{Et si le ratio de Sharpe réel est plus faible que vous ne le pensez~?}{Alors le multiple réel $c$ est plus élevé que ce que l'on pense, et la probabilité d'un drawdown est plus forte. Un $c$ de 0.376 calculé sur un ratio de Sharpe surestimé de 50~\% vaut en réalité 0.564, et la probabilité de perdre la moitié monte à environ 17~\%.}
\end{solution}

\begin{solution}{iq:iv2:betting-and-kelly:28}
Miser environ 19.6~\%, soit $2E[p] - 1$. La loi a posteriori est $\mathrm{Beta}(61, 41)$, de moyenne environ 0.598. L'incertitude sur $p$ ne réduit pas la mise pour un pari unique, car l'espérance du logarithme de la richesse ne dépend de $p$ qu'à travers sa moyenne.

L'espérance du logarithme de la richesse pour une mise $f$ vaut $E[p]\ln(1 + f) + (1 - E[p])\ln(1 - f)$, qui est linéaire en $p$. Seule la moyenne compte donc. L'écart type a posteriori, d'environ 0.048, n'intervient pas.

Miser à la borne basse, $p \approx 0.50$, revient à ne rien miser. L'argument contre~: selon le modèle énoncé, on renonce à toute la croissance espérée, environ 0.019, sans raison. Le modèle intègre déjà l'incertitude.

Les arguments en faveur d'une réduction de la mise concernent le modèle, et non l'arithmétique. Le premier est qu'un a priori uniforme est trop généreux~: les vrais avantages sont rares et petits. Avec un a priori $\mathrm{Beta}(100, 100)$, la moyenne a posteriori vaut $160/300 \approx 0.533$, ce qui donne une mise d'environ 6.7~\%. Le deuxième est la sélection~: il s'agit d'un parieur qu'on a remarqué à cause d'une bonne série. Le troisième est que l'avantage peut changer. Il faut donc fixer l'a priori, puis utiliser la moyenne.

\iqwrong{Diviser par une variance de l'estimateur pour «~pénaliser l'incertitude~». Cela tient dans certains cadres, mais n'a aucun fondement ici. La réduction de l'estimation doit venir d'un meilleur a priori.}
\end{solution}

\begin{solution}{iq:iv2:betting-and-kelly:29}
La probabilité maximale est 64.8~\%, et l'on mise d'abord 0.50~\$. Miser 20~\% à chaque fois donne 0~\%, car trois gains n'amènent qu'à 1.728~\$. Miser tout à chaque fois donne 21.6~\%.

Raisonner à rebours depuis la fin. Avec un pari restant~: à partir de 2~\$ ou plus, l'objectif est atteint~; de 1~\$ à 2~\$, on peut miser assez pour atteindre 2~\$ avec une probabilité de 0.6, et en dessous de 1~\$ la probabilité est nulle. Avec deux paris restants~: à partir de 1.50~\$, on mise pour atteindre 2~\$, et si ce pari est perdu, on se trouve à 1~\$ ou plus avec un pari restant, de sorte que la probabilité vaut $0.6 + 0.4 \times 0.6 = 0.84$. De 1~\$ à 1.50~\$, la probabilité vaut 0.6. De 0.50~\$ à 1~\$, elle vaut $0.6 \times 0.6 = 0.36$.

Avec trois paris restants à 1~\$, miser 0.50~\$ donne $0.6 \times 0.84 + 0.4 \times 0.36 = 0.648$. Les autres mises ne font pas mieux. Il s'agit de programmation dynamique\index{programmation dynamique}, et l'on voit que la fonction valeur est un escalier, de sorte que les seules mises qui méritent d'être considérées sont celles qui tombent sur une marche.

\iqwrong{Utiliser la mise de Kelly parce que «~Kelly est optimal~». Elle maximise le logarithme de la richesse finale, un objectif qui n'est pas celui qui a été fixé.}
\iqfollow{Et si votre objectif était une fonction régulière de la richesse finale, et non une cible stricte~?}{L'escalier disparaît, et la mise varie de façon régulière. Une cible stricte est ce qui produit ces sauts, et aussi ce qui pousse à doubler la mise quand on est en retard.}
\end{solution}

\begin{solution}{iq:iv2:betting-and-kelly:30}
Le rival a l'espérance de richesse la plus élevée, $1.06^{1000}$ contre $1.04^{1000}$. On finit devant avec une probabilité d'environ 95~\%. Les deux sont vrais~: la moyenne du rival provient de trajectoires qui ne se produisent presque jamais.

L'espérance de richesse par pari vaut $1 + 0.2f$, donc une mise plus forte signifie une moyenne plus élevée. Pour comparer les trajectoires, le rapport entre la richesse du joueur et celle du rival après $k$ gains en $n$ paris vaut $(1.2/1.3)^k(0.8/0.7)^{n - k}$. Ce rapport dépasse 1 lorsque $k/n < \ln(0.8/0.7)/\left(\ln(0.8/0.7) + \ln(1.3/1.2)\right) \approx 0.6252$. Avec 1\,000 paris, cela fait au plus 625 gains, et la probabilité binomiale vaut environ 0.95.

Le logarithme typique de la richesse vaut $1000g$~: environ 20.1 pour le joueur et environ 14.7 pour le rival. L'espérance du logarithme de la richesse vaut $1000\ln 1.04 \approx 39.2$ et $1000\ln 1.06 \approx 58.3$. Chacune de ces espérances est très supérieure à ce que l'un ou l'autre joueur observe typiquement. Elle est portée par les trajectoires où la part de gains est élevée, et le rival y gagne.

\iqwrong{Dire que Kelly maximise l'espérance de richesse. Kelly sacrifie la moyenne, et en échange il bat toute autre fraction fixe avec une probabilité qui tend vers 1 quand le nombre de paris croît.}
\end{solution}

\begin{solution}{iq:iv2:betting-and-kelly:31}
La fraction optimale pour la croissance vaut 50~\%. Le rapport moyenne sur variance donne environ 89~\%, et la croissance à cette taille vaut environ $-0.039$ par pari, de sorte qu'on perd de l'argent au fil du temps sur un pari à avantage positif.

Maximiser $0.9\ln(1 + 0.25f) + 0.1\ln(1 - f)$. La dérivée s'annule lorsque $0.225(1 - f) = 0.1(1 + 0.25f)$, donc $0.125 = 0.25f$, et $f = 0.5$. En ce point, la croissance vaut environ 0.0367.

La moyenne vaut $0.9 \times 0.25 - 0.1 = 0.125$. La variance vaut $0.15625 - 0.125^2 \approx 0.1406$. Le rapport vaut 0.889, et avec le moment d'ordre deux brut au dénominateur, il vaut 0.8. Pour une taille de 0.889, la croissance vaut $0.9\ln 1.2222 + 0.1\ln 0.1111 \approx -0.039$.

L'approximation $\mu/\sigma^2$ provient d'un développement à l'ordre deux, qui convient pour de petits paris symétriques. Elle ignore l'asymétrie. Ici une perte de la mise entière est un mouvement important, et le logarithme la sanctionne plus lourdement que le développement.

Dimensionner une stratégie qui vend des options d'après son ratio de Sharpe surestime sa taille appropriée, car le ratio de Sharpe ne voit pas ce que fait le krach. Il faut dimensionner à partir d'une perte de stress~: se demander ce qu'il advient de la richesse dans l'état de perte, et prendre le logarithme de cette richesse.

\iqwrong{Prendre un bon ratio de Sharpe pour une preuve suffisante. Le ratio de Sharpe de cette opération est de 0.33 par pari, et c'est un mauvais candidat pour un dimensionnement par la formule.}
\end{solution}

\begin{solution}{iq:iv2:betting-and-kelly:32}
Avant chaque carte, on mise la fraction $|R - B|/(R + B)$ sur la couleur qui reste le plus nombreuse dans le jeu. On termine avec $2^{52}/\binom{52}{26} \approx \$9.08$ quel que soit l'ordre des cartes, sans aucun risque. Aucune stratégie ne peut donner une richesse finale espérée plus élevée.

Miser de cette façon revient à répartir sa richesse en proportion des chances~: $R/(R + B)$ sur rouge et $B/(R + B)$ sur noir. Si une carte rouge sort, la richesse est multipliée par $2R/(R + B)$. Sur l'ensemble du jeu, les numérateurs donnent $26!\,26!$ quelle que soit la suite des couleurs, et les dénominateurs donnent $52!$. Le multiplicateur vaut donc $2^{52}\,26!\,26!/52!$.

Pour la borne, remarquons qu'un pari à parité déplace la richesse d'une suite à une autre sans en créer. Donc, en sommant sur les $2^{52}$ suites de couleurs, la richesse finale totalise $2^{52}$, quelle que soit la stratégie. Seules $\binom{52}{26}$ suites peuvent se produire, et la richesse n'étant pas négative, leur moyenne est au plus $2^{52}/\binom{52}{26}$. La stratégie proportionnelle l'atteint avec certitude, et par l'inégalité de Jensen aucune stratégie ne donne non plus une espérance du logarithme plus élevée.

\iqwrong{Attendre la dernière carte, où l'on connaît la couleur avec certitude. Cela double l'argent une seule fois. Le pari proportionnel gagne à chaque carte déséquilibrée et n'a aucune variance.}
\iqwhy{Elle illustre la maximisation de la croissance dans sa forme la plus nette~: quand l'issue d'un pari ne laisse plus que peu de possibilités, on est payé en information.}
\end{solution}

\begin{solution}{iq:iv2:betting-and-kelly:33}
L'approximation normale demande environ 1\,075 paris avec une mise de Kelly et environ 476 avec une mise de demi-Kelly. La loi binomiale exacte, en prenant le plus petit nombre à partir duquel la probabilité reste supérieure ou égale à 95~\%, donne 1\,104 et 495.

Pour être en gain avec une confiance de 95~\%, il faut $n g \geq 1.645\sqrt n\,s$, donc $n \geq (1.645\,s/g)^2$, où $g$ est la croissance par pari et $s$ l'écart type du logarithme de la richesse. Pour $f = 0.10$~: $g \approx 0.005008$ et $s \approx 0.0998$. Pour $f = 0.05$~: $g \approx 0.003753$ et $s \approx 0.0498$.

Diviser la mise par deux réduit la croissance d'un quart, mais réduit aussi de moitié son écart type. Donc le rapport $g/s$ augmente de moitié, et le nombre de paris baisse de plus de moitié. Pour le décompte exact, on est en gain lorsque la fréquence des gains dépasse la fraction de paris pour laquelle le logarithme de la richesse reste plat~: $\ln(1/(1 - f))/\ln((1 + f)/(1 - f))$. On prend alors la queue de la loi binomiale. La probabilité croît en dents de scie avec $n$, car les gains sont des nombres entiers, et c'est pour cela que le décompte exact est un peu plus élevé.

\iqwrong{Supposer que le Kelly plein atteint le plus vite n'importe quel objectif. Il maximise la croissance à long terme, mais un objectif exprimé en termes de confiance favorise une mise plus petite.}
\iqfollow{Quelle mise permet d'atteindre le plus vite 95~\%~?}{Lorsque la mise tend vers zéro, $g/s$ tend vers $(p - q)/(2\sqrt{pq})$ et le décompte tombe à environ 268, soit le même que le décompte pour que la somme du pari soit positive. Une mise plus petite est donc meilleure sur cet objectif, mais on paie en richesse atteinte.}
\end{solution}

\begin{solution}{iq:iv2:betting-and-kelly:34}
La richesse finale vaut $2^n h!\,t!/(n + 1)! = 2^n/\left((n + 1)\binom{n}{h}\right)$, avec $h$ piles et $t$ faces, quel que soit l'ordre de leur apparition. Avec 7 piles sur 10 lancers, on termine avec environ 0.776~\$. Quelqu'un qui connaîtrait $p = 0.7$ et jouerait le Kelly termine avec environ 2.28~\$. Avec 10 piles sur 10, on termine avec environ 93.09~\$, et celui qui connaît $p = 0.7$ termine avec $1.4^{10}$, soit environ 28.93~\$~; sur ce parcours, on fait donc mieux.

Une mise de $|2m - 1|$ sur le côté le plus vraisemblable revient à répartir la richesse en $m$ sur pile et $1 - m$ sur face. Ainsi un pile multiplie la richesse par $2m$, et une face par $2(1 - m)$. Avec un a priori uniforme, après $a$ piles et $b$ faces, la moyenne a posteriori vaut $m = (a + 1)/(a + b + 2)$, connue sous le nom de règle de succession de Laplace\index{règle de succession}.

Multiplions maintenant le long d'un chemin quelconque. Les dénominateurs sont $2, 3, \ldots, n + 1$, qui donnent $(n + 1)!$, quel que soit l'ordre des tirages. Chaque pile fournit le suivant de $1, 2, \ldots, h$, et chaque face le suivant de $1, 2, \ldots, t$. Le produit vaut donc $2^n h!\,t!/(n + 1)!$, et il ne dépend que de $h$.

Il existe une vision plus claire. Imaginons qu'on répartisse son dollar entre toutes les valeurs possibles de $p$, en proportion de l'a priori, chaque part jouant le Kelly avec son propre $p$. Le total après $n$ lancers vaut $\int_0^1 2^n p^h(1 - p)^t\,dp$, qui est le même nombre. Le pari bayésien est donc un portefeuille de joueurs de Kelly, et le meilleur d'entre eux domine à long terme.

Comparons avec le cas où $p = 0.7$ est connu~: $2^{10} \times 0.7^7 \times 0.3^3 \approx 2.277$. C'est 2.9 fois ce que l'on a obtenu. En général, la richesse reste en deçà de la meilleure valeur fixe de $p$ d'un facteur au plus égal à $n + 1$, de sorte que le coût de l'apprentissage ne croît que polynomialement, tandis que la richesse procurée par un avantage réel croît exponentiellement.

\iqwrong{Penser que 7 piles sur 10 doivent rapporter de l'argent. Les premiers paris sont faits avec peu de connaissances, et avec si peu de lancers, l'apprentissage coûte plus que ce que rapporte l'avantage.}
\iqwhy{Elle donne un chiffre au coût de ne pas connaître son avantage. Un avantage qu'il faut apprendre vaut moins, et il coûte le plus au début.}
\end{solution}

\begin{solution}{iq:iv2:betting-and-kelly:35}
La probabilité la plus élevée vaut $\Phi(1) \approx 84.1~\%$. On l'atteint en répliquant une option digitale qui rapporte 2~\$ si l'actif finit au-dessus d'un niveau choisi de sorte que l'option coûte exactement 1~\$. Le levier de Kelly donne environ 42.3~\%, et le demi-Kelly environ 26.2~\%. Si l'actif a mal performé en fin d'année, on détient un levier très élevé, puisqu'on n'a plus rien à perdre.

Le marché est complet, donc tout paiement coûte son espérance risque-neutre. Obtenir 2~\$ sur un événement $A$ et rien ailleurs coûte $2Q(A)$, donc il faut $Q(A) \leq \tfrac12$. On veut maximiser $P(A)$ pour ce $Q(A)$. Par le lemme de Neyman-Pearson\index{lemme de Neyman-Pearson}, on prend les états où $dP/dQ$ est le plus grand, c'est-à-dire ceux où l'actif a le plus progressé.

Soit $B$ un mouvement brownien sous $P$. Sous $Q$, $B^Q_t = B_t + St$ est un mouvement brownien. Donc $A = \{B^Q_1 \geq 0\}$ vérifie $Q(A) = \tfrac12$, et $P(A) = P(B_1 + 1 \geq 0) = \Phi(1)$. Plus généralement, avec une richesse $w$, un objectif $G$ et un temps $\tau$ restant, la réponse est $\Phi\left(\Phi^{-1}(w/G) + S\sqrt\tau\right)$. La volatilité de l'actif n'apparaît pas, seul le ratio de Sharpe intervient.

La richesse à l'instant $t$ est la valeur de l'option digitale, $W_t = 2\Phi\left(B^Q_t/\sqrt{1 - t}\right)$, et les dollars de risque pris par unité de volatilité de l'actif valent $\partial W/\partial B^Q = 2\varphi\left(B^Q_t/\sqrt{1 - t}\right)/\sqrt{1 - t}$. Exprimé en multiple de la richesse, cela vaut $\varphi(z)/\left(\Phi(z)\sqrt{1 - t}\right)$ avec $z = B^Q_t/\sqrt{1 - t}$. Il croît à mesure que $z$ diminue et que l'échéance approche, sans borne. Donc si l'on est en retard en fin d'année, on détient un levier très élevé.

Pour comparaison, le Kelly donne $\ln W_1 = \tfrac12 + B_1$, et il faut $B_1 \geq \ln 2 - \tfrac12 \approx 0.193$, soit une probabilité d'environ 42.3~\%. Le demi-Kelly donne $\ln W_1 = 0.375 + \tfrac12 B_1$, qui exige $B_1 \geq 0.636$, soit une probabilité d'environ 26.2~\%.

\iqwrong{Croire qu'une probabilité de succès plus élevée est forcément bonne. Les 15.9~\% de trajectoires perdantes finissent à zéro, un résultat que ni le Kelly ni le demi-Kelly ne peuvent donner, et la meilleure stratégie utilise un levier illimité juste avant l'échéance.}
\iqwhy{C'est ce vers quoi est poussé un gérant qui fait face à un objectif difficile avec une option gratuite à la hausse, et la raison pour laquelle on limite le levier en fin de période.}
\end{solution}
